\ifdefined\HMTAPPFormal
\chapter{APP: soporte, operaciones y caminos}
\label{chap:app}
\index{APP}

\noindent
La sigla APP procede de \emph{All Possible Paths}. Designa una única
estructura de caminos sobre un soporte aritmético finito: el grafo
subyacente tiene ochenta y un vértices, mientras que el conjunto de caminos
finitos, reunido sobre todas las longitudes, es numerable. Su matriz
aritmética principal es la tabla multiplicativa nonádica. La suma es la ley
interna cuya acumulación genera cada fila y cada columna de esa tabla y actúa
además como observable de comparación sobre el mismo soporte. APP reúne así
suma, producto y
composición de trayectorias antes de efectuar la proyección ternaria o la
evaluación arquimediana. La definición siguiente fija de manera exhaustiva
el universo de caminos considerado.

\section{Dominio, generadores y relaciones preformales del soporte APP}

El soporte de APP es un toro aritmético finito: cada una de sus dos
coordenadas recorre las nueve clases residuales y vuelve a su origen después
de nueve traslaciones unitarias. Sus ochenta y un puntos son, por tanto,
pares de clases módulo nueve. El grafo de Cayley asociado conserva esta
estructura cíclica: un vértice representa una posición, una arista orientada
representa un desplazamiento cardinal unitario y una sucesión finita de
aristas representa un camino. De este modo, \emph{todos los caminos
posibles} designa el conjunto de las palabras cardinales finitas con punto
inicial en el toro.

Cada vértice recibe simultáneamente dos evaluaciones aritméticas. La
\emph{hoja aditiva} publica la suma de sus dos coordenadas y la
\emph{hoja multiplicativa} publica su producto. El término \emph{hoja}
designa aquí una evaluación del mismo punto y del mismo camino; APP conserva
un único soporte y dos funciones sobre él.

La diferencia operativa entre ambas hojas es exacta. En una fila o columna,
la suma actúa por traslación y acumula reiteradamente el mismo incremento.
La multiplicación por una unidad también permuta las clases; la
multiplicación por un divisor de cero pliega varias clases sobre un mismo
residuo. Para conservar la información que ese plegado oculta, APP eleva
cada evaluación a sus coordenadas residuo--cociente:
\[
 i+j=R^+_{ij}+9q^+_{ij},
 \qquad
 ij=R^\times_{ij}+9q^\times_{ij}.
\]
Los pares \((R^+,q^+)\) y \((R^\times,q^\times)\) recuperan los valores
enteros anteriores a la reducción. Al concatenar bloques, el acarreo aditivo
\(\kappa_9\) registra la discrepancia entre sumar representantes enteros y
sumar primero sus residuos; junto con ambos cocientes, hace composicional la
recuperación. Así, la acumulación aditiva construye las trayectorias visibles,
mientras el levantamiento residuo--cociente despliega las fibras creadas por el
plegado multiplicativo.

Para un camino \(\gamma=(x_0,\ldots,x_r)\), la genealogía visible es la
sucesión de posiciones y pares de evaluaciones
\[
 \mathcal G_{\rm vis}(\gamma)
 =\bigl(\gamma;(\Sigma(x_k),\Pi(x_k))_{k=0}^{r}\bigr).
\]
Su genealogía enriquecida añade, en cada paso, los cocientes, el acarreo y
la orientación transportada. La proyección enriquecida--visible conserva el
camino y los residuos y olvida precisamente esas coordenadas de memoria. Esta
distinción será la entrada del estado completo del TRIT y del Topological Phase Kernel.
El dominio completo de (104\,976) condiciones iniciales de doble cursor y
las (468) fibras de emisión que producen se publican íntegramente en el
\cref{app:catalogo-semillas-app-rev9}; allí se distingue cada semilla de la
emisión a la que pertenece y se conserva la multiplicidad de la fibra.

\section{Soporte aritmético, traslaciones y observables}

Sea
\[
\dr(n)=1+((n-1)\bmod9)\in\{1,\ldots,9\}.
\]
Esta aplicación escoge el representante \(9\) para la clase nula. Definimos
\(X_9=(\ZZ/9\ZZ)^2\), sobre el que actúa el grupo de traslaciones
\((\ZZ/9\ZZ)^2\), y fijamos el conjunto simétrico de generadores
\[
 \mathscr D=\{(1,0),(0,1),(-1,0),(0,-1)\}.
\]
El grafo de Cayley
\[
 \GraphAPP=\operatorname{Cay}(X_9,\mathscr D)
\]
posee \(81\) vértices y cuatro aristas orientadas con origen en cada vértice.
Un camino orientado de longitud \(r\) es una sucesión
\(\gamma=(x_0,\ldots,x_r)\) en \(X_9\) tal que
\(x_{k+1}-x_k\in\mathscr D\) para todo \(k<r\). Denotamos por
\(\operatorname{Path}(\GraphAPP)\) el conjunto de todos los caminos finitos
así definidos.

\begin{definition}[Estructura APP]
La estructura aritmética de caminos posibles es la tupla
\[
 \AAPP=(X_9,\GraphAPP,\Sigma,\Pi,
        \operatorname{Path}(\GraphAPP)),
\]
donde los dos observables de vértice son
\[
\Sigma(i,j)=\dr(i+j),
\qquad
\Pi(i,j)=\dr(ij),
\qquad1\le i,j\le9.
\]
\end{definition}

\subsection*{Marco cardinal y palabras de trayectoria}

Para hacer explícita la orientación del toro, usamos coordenadas
\((\text{fila},\text{columna})\), con la fila creciente hacia el sur y la
columna creciente hacia el este, y escribimos
\[
 \delta(N)=(-1,0),\qquad \delta(E)=(0,1),\qquad
 \delta(S)=(1,0),\qquad \delta(O)=(0,-1).
\]
Sea \(\mathscr D_{\rm card}=\{E,N,O,S\}\), donde \(O\) denota el oeste. Dado un punto inicial
\(x_0\in X_9\) y una palabra \(w=d_1\cdots d_r\in
\mathscr D_{\rm card}^{r}\), su trayectoria es
\[
 \gamma(x_0;w)_k
 =x_0+\sum_{j=1}^{k}\delta(d_j)\pmod 9,
 \qquad 0\leq k\leq r.
\]
La correspondencia \((x_0,w)\mapsto\gamma(x_0;w)\) es biyectiva entre
\(X_9\times\mathscr D_{\rm card}^{r}\) y los caminos orientados de
longitud \(r\). En particular, existen exactamente \(81\cdot4^r\) caminos
de esa longitud cuando se permiten retornos y retrocesos.

La palabra \(w\) cierra en el toro precisamente cuando
\[
\#S(w)-\#N(w)\equiv0\pmod9,
\qquad
\#E(w)-\#O(w)\equiv0\pmod9.
\]
Para una palabra cerrada, el desplazamiento entero calculado antes de reducir
módulo nueve define el vector de enrollamiento
\[
 \operatorname{wind}(w)=\frac1{9}
 \bigl(\#S(w)-\#N(w),\#E(w)-\#O(w)\bigr)\in\ZZ^2.
\]
La inversión de la orientación revierte el orden de la palabra, intercambia
\(E\leftrightarrow O\) y \(N\leftrightarrow S\), y transforma
\(\operatorname{wind}(w)\) en su opuesto. Este marco cardinal es el que se
transporta posteriormente en el núcleo de fase y forma parte de la definición
de cada trayectoria.

Las expresiones se escriben aquí en la carta de representantes
\(\{1,\ldots,9\}\). La fórmula
\(\widetilde\Sigma(i,j)=\dr(i+j-1)\) es una carta trasladada. Representa el
mismo toro sólo cuando la traslación se aplica también a índices, condiciones
iniciales, trayectorias posicionales y etiquetas. Todas las demostraciones siguientes emplean
\(\Sigma(i,j)=\dr(i+j)\) y no alternan ambas cartas.

\begin{figure}[H]
\centering
\resizebox{\textwidth}{!}{\begin{tikzpicture}[font=\scriptsize,cell/.style={minimum width=4.2mm,minimum height=4.2mm,inner sep=0pt,draw=hmtgray!35}]
  \begin{scope}[xshift=0cm]
    \node[font=\bfseries\scriptsize,hmtblue,align=center] at (1.72,.92)
      {matriz aritmética APP\\[-.5mm]\(\Pi(i,j)=\dr(ij)\)};
    \foreach \i in {1,...,9}{
      \node[font=\bfseries] at (-.35,{-(\i-1)*.43}) {\i};
      \foreach \j in {1,...,9}{
        \pgfmathtruncatemacro{\v}{mod(\i*\j-1,9)+1}
        \pgfmathtruncatemacro{\border}{(\i==9)||(\j==9)}
        \pgfmathtruncatemacro{\rad}{(mod(\i,3)==0)||(mod(\j,3)==0)}
        \ifnum\border=1
          \node[cell,fill=hmtlightblue,text=hmtblue,font=\bfseries] at ({(\j-1)*.43},{-(\i-1)*.43}) {\v};
        \else
          \ifnum\rad=1
            \node[cell,fill=hmtlightviolet] at ({(\j-1)*.43},{-(\i-1)*.43}) {\v};
          \else
            \node[cell,fill=white] at ({(\j-1)*.43},{-(\i-1)*.43}) {\v};
          \fi
        \fi
      }
    }
    \foreach \j in {1,...,9}{\node[font=\bfseries] at ({(\j-1)*.43},.38) {\j};}
    \draw[hmtgold,very thick,rounded corners=1pt] (1.075,-1.075) rectangle (1.935,-1.935);
    \draw[hmtblue,very thick] (3.72,.22)--(3.72,-3.72)--(-.22,-3.72);
    \node[hmtblue,anchor=north] at (1.75,-4.02) {subconjunto coordenado de la clase nula};
  \end{scope}

  \begin{scope}[xshift=5.25cm]
    \node[draw=hmtblue,rounded corners=2pt,fill=hmtlightblue,
      text width=6.25cm,align=left,inner sep=8pt] at (2.9,-1.35) {%
      \textbf{Ley interna de acumulación}\par\smallskip
      \[
        \Pi(i,j)=\dr(\underbrace{i+\cdots+i}_{j\ {\rm sumandos}})
        =\dr(ij).
      \]
      La suma genera las filas multiplicativas y, en la dirección
      transversal, las columnas. El observable
      \(\Sigma(i,j)=\dr(i+j)\) conserva esa ley elemental para comparar
      ambas lecturas sobre el mismo punto.};
    \node[draw=hmtgold,rounded corners=2pt,fill=hmtlightgold,
      text width=6.25cm,align=center,inner sep=8pt] at (2.9,-3.65) {%
      \(\displaystyle
      \AAPP=(X_9,\GraphAPP,\Pi;\Sigma,
      \operatorname{Path}(\GraphAPP))\)\par
      una sola APP, una matriz principal, una ley aditiva interna y cuatro
      direcciones orientadas \(N,E,S,O\).};
  \end{scope}
\end{tikzpicture}
}
\caption{La única matriz APP y su ley interna. La tabla multiplicativa
nonádica ocupa el cuerpo principal; cada una de sus filas se genera por
acumulación aditiva. El sombreado violeta señala las filas o columnas no
unitarias; la línea azul identifica el subconjunto determinado por el
representante \(9\) de la clase nula; el marco dorado destaca el bloque
central \(2\times2\).}
\label{fig:app-ortograma}
\end{figure}

Cada punto de \(X_9\) recibe simultáneamente los valores de producto y suma,
pero no pertenece por ello a dos matrices APP distintas. Los caminos de
\(\operatorname{Path}(\GraphAPP)\) recorren un único soporte; \(\Pi\)
publica su tabla principal y \(\Sigma\) conserva la operación interna que
genera las acumulaciones y permite medir su incompatibilidad.

\begin{proposition}[Estructura de las filas y columnas]
\label{prop:app-latina-unidades}
La matriz de \(\Sigma\) es un cuadrado latino de orden nueve. La matriz de
\(\Pi\) no lo es; sin embargo, para cada
\(u\in(\ZZ/9\ZZ)^\times=\{1,2,4,5,7,8\}\), la fila y la columna
correspondientes a \(u\) son permutaciones de las nueve clases residuales.
\end{proposition}

\begin{proof}
Fijado \(i\), la aplicación \(j\mapsto i+j\) es una traslación de
\(\ZZ/9\ZZ\); lo mismo vale al fijar \(j\). Por ello cada representante
aparece exactamente una vez en cada fila y columna de \(\Sigma\). Para
\(\Pi\), la multiplicación por \(u\) permuta las clases si y sólo si \(u\)
es una unidad. Las filas correspondientes a \(3,6,9\) poseen repeticiones,
por lo que la propiedad latina no se extiende a toda la matriz
multiplicativa.
\end{proof}

\section{Estructura multiplicativa del anillo local}

Cada fila de la tabla multiplicativa es una suma repetida. La fila \(i\)
acumula \(i\) a lo largo de una dirección; la columna \(j\) acumula \(j\) a
lo largo de la dirección perpendicular. La celda \((i,j)\) es la intersección
de ambas acumulaciones y devuelve \(ij=ji\). La tabla representa así producto
y conmutatividad mediante dos acumulaciones transversales sobre la misma
retícula.

La segunda fila es
\[
2,4,6,8,1,3,5,7,9.
\]
Antes del retorno modular aparecen los pares y después los impares. La tercera
fila es
\[
3,6,9,3,6,9,3,6,9,
\]
y hace visible el radical. Sea
\[
 R:=\ZZ/9\ZZ,
 \qquad
 \mathfrak m:=(3)=\{\overline0,\overline3,\overline6\}.
\]
La multiplicación por tres define
\[
 M_3:R\longrightarrow R,
 \qquad x\longmapsto3x,
 \qquad
 \ker M_3=\operatorname{im}M_3=\mathfrak m.
\]
Por el primer teorema de isomorfía induce una identificación lineal
\begin{equation}
 \overline M_3:R/\mathfrak m\xrightarrow{\ \sim\ }\mathfrak m,
 \qquad \overline x\longmapsto3x,
 \label{eq:multiplicador-tres-radical-rev3}
\end{equation}
de espacios vectoriales sobre \(\mathbb F_3\). No es un isomorfismo de
anillos: \(\mathfrak m^2=0\), mientras el cociente es un cuerpo.

En representantes positivos, las tres fases de
\eqref{eq:multiplicador-tres-radical-rev3} son
\begin{equation}
 369\longmapsto693\longmapsto936\longmapsto369,
 \label{eq:orbita-radical-369-rev3}
\end{equation}
producidas por el desplazamiento de fase \(j\mapsto j+1\). La fila
correspondiente a \(6\equiv-3\pmod9\) invierte la orientación del mismo
ciclo. Esta órbita radical no debe confundirse con el entero
\(369\,693\,100\), que aparecerá después como número de cifras de
\(9^{9^9}\).

En reducción directa módulo tres,
\[
3,6,9\longmapsto0,0,0.
\]
Si primero extraemos el factor tres y guardamos su coordenada interna,
\[
\eta_{\mathfrak m}(3a)=a\bmod3,
\qquad
3,6,9\longmapsto1,2,0.
\]
La segunda aplicación conserva la coordenada interna del ideal generado por
tres. No debe confundirse con la proyección ternaria orientada que se definirá
en la sección dedicada al TRIT\@.

\section{Subconjunto coordenado de la clase nula}

En la gráfica del observable multiplicativo,
\[
\Pi(9,j)=\Pi(i,9)=9.
\]
La última fila y la última columna forman una unión de diecisiete posiciones,
todas con valor nueve. No constituyen una frontera intrínseca del toro
aritmético: son el subconjunto de la carta en el que al menos una coordenada
representa la clase nula mediante \(9\). Existen además cuatro representantes nulos interiores en
\[
(3,3),(3,6),(6,3),(6,6).
\]
La figura distingue así este subconjunto y los ceros internos del ideal
no unitario. El símbolo adicional empleado más adelante en la completación
cúbica pertenece a un espacio ampliado y no se identifica con ninguno de
ellos.

\section{Extensión residuo--cociente}

La reducción módulo nueve conserva el residuo y elimina el cociente de la
división euclídea. Para estudiar ambas componentes definimos, en cada punto,
\[
i+j=9q^+_{ij}+R^+_{ij},
\qquad
ij=9q^\times_{ij}+R^\times_{ij},
\]
donde \(R^+=\Sigma\), \(R^\times=\Pi\) y los \(q\) conservan el cociente.
La extensión residuo--cociente de APP se denota por
\[
\AAPPlift=(R^+,q^+;R^\times,q^\times).
\]

Sea \(s_9:\ZZ/9\ZZ\to\{1,\ldots,9\}\) la sección de representantes fijada
por \(\dr\). El acarreo aditivo asociado a esta sección es
\[
 \kappa_9(a,b)
 =\frac{s_9(a)+s_9(b)-s_9(a+b)}9\in\ZZ.
\]

\begin{lemma}[Identidad cocíclica del acarreo]
\label{lem:app-cociclo-acarreo}
Para cualesquiera \(a,b,c\in\ZZ/9\ZZ\),
\[
 \kappa_9(a,b)+\kappa_9(a+b,c)
 =\kappa_9(b,c)+\kappa_9(a,b+c).
\]
\end{lemma}

\begin{proof}
Ambos miembros son iguales a
\[
 \frac{s_9(a)+s_9(b)+s_9(c)-s_9(a+b+c)}9.\qedhere
\]
\end{proof}

Este cociclo registra la discrepancia entre sumar representantes enteros y
sumar primero en el cociente. Es la forma composicional de la coordenada de
cociente que reaparecerá en la concatenación base \(1000\).

Los totales brutos son
\[
\sum_{i,j}(i+j)
=9\sum_i i+9\sum_j j
=2\cdot9\cdot45
=810,
\]
y
\[
\sum_{i,j}ij
=\left(\sum_i i\right)\left(\sum_j j\right)
=45^2
=2025.
\]
Las sumas visibles se obtienen de las tablas:
\[
\sum R^+=405,
\qquad
\sum R^\times=459.
\]
Por tanto, las sumas totales de las coordenadas de cociente son
\[
\sum q^+=\frac{810-405}{9}=45,
\qquad
\sum q^\times=\frac{2025-459}{9}=174.
\]

\begin{exactbox}[title={Descomposición exacta de la diferencia integrada}]
\[
\boxed{
2025-810
=1215
=(459-405)+9(174-45)
=54+9\cdot129.}
\]
La diferencia entre las sumas de residuos es \(54\), mientras que la
diferencia entre las sumas de cocientes es \(129\). La igualdad recompone la
diferencia total entre los observables enteros y demuestra que la componente
de cociente contiene información aritmética no determinada por la clase
residual.
\end{exactbox}

\fi

\ifdefined\HMTAPPDevelopments
\chapter{Operadores aritméticos y geometría local del espacio de caminos}
\label{chap:app-operadores}

\section{Localización del exceso en el radical nilpotente}

El grupo de unidades de \(\ZZ/9\ZZ\) es
\[
(\ZZ/9\ZZ)^\times=\{1,2,4,5,7,8\},
\]
y el ideal máximo es
\[
\mathfrak m=(3)=\{0,3,6\}=\{9,3,6\}_{\rm visual},
\qquad\mathfrak m^2=0.
\]
Cada fila unitaria de \(\Pi\) es una permutación de \(1,\ldots,9\) y suma
\(45\). Las filas \(3\) y \(6\) suman \(54\); la fila \(9\), \(81\).
Luego
\[
459=6\cdot45+2\cdot54+81
\]
y
\[
459-405=2(54-45)+(81-45)=54.
\]
Todo el exceso se localiza en el sector no unitario.

La unión de filas y columnas \(3,6,9\) contiene \(45\) celdas y suma
\(297\). Su cuadrado central \(3\times3\) suma \(81=3\cdot27\); los brazos
exteriores, \(216=3\cdot72\). Así,
\[
297=216+81=3(72+27)=3\cdot99.
\]
La interpretación de \(72\) y \(27\) se obtiene más adelante a partir del
bloque central, seleccionado de manera única.

\section{Esperanza condicional respecto de la proyección \(\mathbb Z/9\mathbb Z\to\mathbb Z/3\mathbb Z\)}

La reducción canónica
\(\ZZ/9\ZZ\to\ZZ/3\ZZ\) determina la partición
\[
 C_1=\{1,4,7\},\qquad
 C_2=\{2,5,8\},\qquad
 C_0=\{3,6,9\}.
\]
La esperanza condicional uniforme promedia cada observable sobre los nueve puntos de cada bloque
\(C_a\times C_b\). Las matrices agregadas son
\[
 M_{\Pi}=
 \begin{pmatrix}4&5&6\\5&4&6\\6&6&9\end{pmatrix},
 \qquad
 M_{\Sigma}=
 \begin{pmatrix}5&6&4\\6&4&5\\4&5&6\end{pmatrix}.
\]
Puesto que cada bloque contiene nueve posiciones,
\[
 9\sum_{a,b}(M_\Pi)_{ab}=459,\qquad
 9\sum_{a,b}(M_\Sigma)_{ab}=405.
\]
Así, el promediado conserva los totales de ambos observables y localiza de
nuevo su diferencia:
\[
 \sum M_\Pi-\sum M_\Sigma=51-45=6,\qquad 9\cdot6=54.
\]

\begin{theorem}[Espectro del operador multiplicativo agregado]
\label{thm:app-espectro-agregado}
El operador matricial \(M_\Pi\) satisface
\[
 \det M_\Pi=-9
\]
y su espectro real es
\[
 \operatorname{Spec}(M_\Pi)
 =\{-1,\,9-6\sqrt2,\,9+6\sqrt2\}.
\]
En particular, la forma cuadrática asociada posee firma \((2,1)\).
\end{theorem}

\begin{proof}
La expansión del determinante da \(-9\). El polinomio característico
factoriza como
\[
 (\lambda+1)\bigl((\lambda-9)^2-72\bigr).
\]
Sus raíces son las indicadas y \(9-6\sqrt2>0\); por tanto existen dos
direcciones positivas y una negativa.
\end{proof}

La descomposición espectral admite una forma integral sobre una subretícula
invariante. Para
\[
 v_-=(1,-1,0),\qquad v_+=(1,1,0),\qquad v_0=(0,0,1),
\]
se tiene \(M_\Pi v_-=-v_-\), mientras que la restricción de \(M_\Pi\) al
submódulo generado por \(v_+,v_0\) está representada por
\[
 \begin{pmatrix}9&6\\12&9\end{pmatrix}
 =3H,\qquad
 H=\begin{pmatrix}3&2\\4&3\end{pmatrix}\in SL(2,\ZZ).
\]
Los autovalores de \(H\) son
\[
 3\pm2\sqrt2=(1\pm\sqrt2)^2.
\]
Los tres vectores \(v_-,v_+,v_0\) generan una subretícula de índice dos en
\(\ZZ^3\), pues el determinante de la matriz que los tiene por columnas vale
\(-2\). Por tanto, no se afirma aquí una suma directa integral de toda
\(\ZZ^3\): sobre esa subretícula invariante, el operador agregado separa una
dirección de autovalor \(-1\) y un submódulo hiperbólico gobernado por la
unidad cuadrática \(3+2\sqrt2\); sobre \(\RR^3\), la descomposición espectral
es completa.

\section{Conmutador de los proyectores aritméticos y separación de sus imágenes}

La diferencia entre \(\Sigma\) y \(\Pi\) no se reduce a sus sumas globales.
Para \(d\ge2\), sea \(\mathscr A_d=\{1,\ldots,9\}^d\) y definamos
\[
 \Sigma_d(x)=\dr(x_1+\cdots+x_d),\qquad
 \Pi_d(x)=\dr(x_1\cdots x_d).
\]
En el espacio de Hilbert \(\ell^2(\mathscr A_d)\), provisto de la medida
uniforme, denotemos por \(P_{\Sigma,d}\) y \(P_{\Pi,d}\) las esperanzas
condicionales ortogonales sobre las funciones medibles respecto de
\(\Sigma_d\) y \(\Pi_d\), respectivamente. La tabla conjunta normalizada es
\[
 C^{(d)}_{ab}
 =\frac{N^{(d)}_{ab}}{\sqrt{n_a m_b}},
\quad
 N^{(d)}_{ab}
 =\#\{x:\Sigma_d(x)=a,\ \Pi_d(x)=b\},
\]
donde \(n_a=\sum_bN^{(d)}_{ab}\) y \(m_b=\sum_aN^{(d)}_{ab}\).

En el soporte bidimensional original, la correlación conjunta exacta es
\[
N^{(2)}=
\begin{pmatrix}
0&0&2&0&0&2&3&0&2\\
3&0&2&0&0&2&0&0&2\\
0&2&0&0&2&0&0&2&3\\
0&0&2&3&0&2&0&0&2\\
0&0&2&3&0&2&0&0&2\\
0&2&0&0&2&0&0&2&3\\
3&0&2&0&0&2&0&0&2\\
0&0&2&0&0&2&3&0&2\\
0&2&0&0&2&0&0&2&3
\end{pmatrix}.
\]
Las filas, ordenadas por el valor de \(\Sigma\), suman nueve; las columnas,
ordenadas por el valor de \(\Pi\), suman
\((6,6,12,6,6,12,6,6,21)\). Así queda determinada por enumeración completa
la correlación conjunta de los dos observables antes de introducir cualquier
aplicación de evaluación posterior.

\begin{proposition}[Conmutadores en dimensiones dos y tres]
\label{prop:app-no-conmutatividad}
Los dos pares de proyectores satisfacen
\[
 \|[P_{\Sigma,2},P_{\Pi,2}]\|=\frac{\sqrt2}{3},
 \qquad
 \|[P_{\Sigma,3},P_{\Pi,3}]\|=\frac{\sqrt3}{4}.
\]
En particular, los observables aditivo y multiplicativo no inducen
descomposiciones simultáneamente compatibles.
\end{proposition}

\begin{proof}
Los valores singulares \(s\) de \(C^{(d)}\) son los cosenos de los ángulos
principales entre los dos subespacios. En el bloque bidimensional
correspondiente, la norma del conmutador es \(s\sqrt{1-s^2}\). La enumeración
exacta de las \(9^d\) palabras produce, para \(d=2\), los autovalores no
triviales
\[
 s^2\in\left\{\frac57,\frac13,\frac13\right\}.
\]
El máximo de \(s^2(1-s^2)\) es \(2/9\), de donde resulta
\(\sqrt2/3\). Para \(d=3\), la misma construcción da el máximo
\(s^2(1-s^2)=3/16\), alcanzado en \(s^2=1/4\), y por tanto la norma es
\(\sqrt3/4\). En ambos casos el conmutador es no nulo.
\end{proof}

\begin{proposition}[Intersección de los subespacios observables]
\label{prop:app-interseccion-proyectores}
La enumeración exacta de las palabras de \(\mathscr A_d\), para
\(2\leq d\leq15\), verifica
\[
 \operatorname{Ran}P_{\Sigma,d}\cap
 \operatorname{Ran}P_{\Pi,d}
 =\CC\mathbf 1.
\]
En ese intervalo de dimensiones, las funciones constantes constituyen el
único subespacio común a las dos descomposiciones aritméticas.
\end{proposition}

\begin{proof}
Para dos proyecciones ortogonales, la dimensión de la intersección de sus
rangos coincide con la multiplicidad del valor singular \(1\) en la matriz
de solapamiento \(C^{(d)}\). El cálculo entero de las tablas
\(N^{(d)}\), seguido del cálculo exacto de esa multiplicidad, da valor uno
para cada \(d=2,\ldots,15\). El vector constante pertenece a ambos rangos;
por tanto genera toda la intersección.
\end{proof}

La sucesión de normas no se extrapola a partir de los dos primeros casos.
En particular, para \(d=4\) el mismo cálculo da
\[
 \|[P_{\Sigma,4},P_{\Pi,4}]\|
 =\sqrt{\frac{2618939}{22240656}}
 \simeq0.3431538652,
\]
que no coincide con la prolongación elemental \(2/(d+1)\). Este caso
constituye un contraejemplo para cualquier fórmula inferida sólo de
\(d=2,3\).

\section{Bloque central y descomposición espectral}

La caracterización del bloque central se obtiene mediante una condición de
congruencia. Para dos representantes consecutivos \(k,k+1\), consideremos
\[
B_k=
\begin{pmatrix}
k^2&k(k+1)\\
k(k+1)&(k+1)^2
\end{pmatrix}\pmod9.
\]

\begin{theorem}[Unicidad del bloque diagonal adyacente]
\label{thm:app-bloque-central}
Existe un único bloque \(B_k\) cuyas entradas diagonales coinciden. Está
determinado por \(k=4\) y es
\[
B_c=
\begin{pmatrix}7&2\\2&7\end{pmatrix}.
\]
\end{theorem}

\begin{proof}
La igualdad de las diagonales equivale a
\[
k^2\equiv(k+1)^2\pmod9,
\]
es decir, \(2k+1\equiv0\pmod9\). Como \(2\) es una unidad de
\(\ZZ/9\ZZ\), la congruencia tiene la única solución \(k\equiv4\pmod9\).
\end{proof}

Por tanto, las posiciones \(4,5\) del observable multiplicativo forman
\[
B_c=
\begin{pmatrix}7&2\\2&7\end{pmatrix}
=7I+2R,
\qquad
R=\begin{pmatrix}0&1\\1&0\end{pmatrix},
\qquad R^2=I.
\]
Con \(P_\pm=(I\pm R)/2\),
\[
B_c=9P_++5P_-.
\]
\begin{theorem}[Jerarquía tipada de la separación suma--producto]
\label{thm:app-jerarquia-cuatro-dos-seis-cincuentaycuatro}
El bloque central tiene salto espectral local (9-5=4). Su coeficiente
fuera de la diagonal es (2). La compresión por las tres clases módulo tres
produce la diferencia agregada (51-45=6), y el despliegue sobre las nueve
posiciones produce el defecto global
\[
 9\cdot6=459-405=54.
\]
Los enteros (4), (2), (6) y (54) pertenecen, respectivamente, al
salto espectral, al acoplamiento local, a la compresión modular y a la
integración bidimensional.
\end{theorem}

\begin{proof}
La descomposición \(B_c=7I+2R=9P_++5P_-\) da el salto y el
coeficiente local. Las matrices agregadas \(M_\Pi,M_\Sigma\) construidas
arriba satisfacen
\(\sum M_\Pi=51\) y \(\sum M_\Sigma=45\). Cada entrada agregada representa
nueve posiciones, por lo que el despliegue multiplica la diferencia por
nueve y recupera los totales (459) y (405).
\end{proof}
Sus concatenaciones ordenadas producen los bloques \(72\) y \(27\):
\[
72+27=99,
\qquad
72-27=45=\det B_c.
\]
La concatenación de las entradas diagonales y antidiagonales produce,
respectivamente, \(77\) y \(22\). Por tanto,
\[
 72+27=77+22=99,\qquad
 77-22=55=54+1.
\]
Estas igualdades expresan dos particiones exactas del mismo bloque; no
identifican por sí solas una constante externa.

\begin{proposition}[Relación local--global del sector radical]
\label{prop:app-local-global-radical}
La partición por filas del bloque central reproduce, con factor de escala
tres, la partición del sector radical:
\[
 3\cdot72=216,\qquad
 3\cdot27=81,\qquad
 3\cdot99=297.
\]
\end{proposition}

\begin{proof}
El bloque \(C_0\times C_0\) contiene nueve representantes nulos y suma
\(81\). Los cuatro brazos restantes de las filas y columnas de \(C_0\)
suman \(216\). Las identidades siguen de \(72+27=99\).
\end{proof}

Además,
\[
M_c:=\frac{B_c}{\sqrt{45}}
=\exp\left(\frac12\log\frac95\,R\right)
\in SO^+(1,1).
\]
Si \(\eta=\operatorname{diag}(1,-1)\), la matriz normalizada \(M_c\) satisface
\(M_c^{\mathsf T}\eta M_c=\eta\), \(\det M_c=1\) y pertenece a
\(SO^+(1,1)\). Esta realización finita se limita al bloque \(2\times2\)
construido; su interpretación espacio-temporal requiere el puente dimensional
desarrollado en el tratado de Mecánica Dimensional.

\begin{proposition}[Familia espectral APP--TRIT y lectura
Markov--Lorentz local]
\label{prop:app-trit-espectral-markov-lorentz-rev8}
Para \(\tau\in\{-1,0,+1\}\), sea
\[
 B_\tau=(4+3\tau)I+(5-3\tau)R.
\]
Entonces
\[
 \operatorname{Spec}(B_\tau)=\{9,\,6\tau-1\},
 \qquad
 B_\tau=9P_++(6\tau-1)P_-,
\]
y el modo diferencial recupera la orientación mediante
\[
 \tau=\frac{\lambda_-+1}{6}.
\]
En particular, \(B_{+1}=B_c\). Su normalización estocástica
\[
 P_B:=\frac{1}{9}B_c=P_++\frac59P_-
\]
satisface, para todo \(n\geq0\),
\[
 P_B^n=P_++\left(\frac59\right)^n P_-.
\]
Por tanto, la contracción del modo antisimétrico es
\(\kappa_B=5/9\) y la brecha espectral de esta cadena local es \(4/9\).
Si
\[
 \eta_c:=\frac12\log\frac95,
\]
la normalización lorentziana del mismo bloque cumple
\[
 \frac{B_c}{\sqrt{45}}=\exp(\eta_c R),
 \qquad
 \boxed{\frac59=e^{-2\eta_c}}.
\]
\end{proposition}

\begin{proof}
Como \(R\) actúa con autovalores \(+1\) y \(-1\) sobre
\(\operatorname{im}P_+\) y \(\operatorname{im}P_-\), respectivamente, los
dos autovalores de \(B_\tau\) son
\[
 (4+3\tau)+(5-3\tau)=9,\qquad
 (4+3\tau)-(5-3\tau)=6\tau-1.
\]
La fórmula de las potencias sigue de
\(P_\pm^2=P_\pm\) y \(P_+P_-=0\). Finalmente,
\(e^{-2\eta_c}=e^{-\log(9/5)}=5/9\).
\end{proof}

Esta coincidencia identifica dos lecturas exactas del bloque local: la
memoria diferencial decrece con el mismo cociente que codifica la rapidez
en la realización descompuesta. No convierte \(B_\tau\) en el núcleo de
Markov global de TPK, ni \(5/9\) en una constante física. Tampoco extiende
por sí sola la realización \(SO^+(1,1)\) a un espacio de Minkowski
\(3+1\): esos pasos requieren mapas dimensionales adicionales.

\section{Curvatura discreta inducida por la interacción de las hojas suma y producto y su defecto de conmutación}

En el toro escribimos, ahora en residuos,
\[
S(x,y)=x+y,
\qquad
P(x,y)=xy\pmod9.
\]
Para toda función \(T:X_9\to\ZZ/9\ZZ\), fijamos las diferencias progresivas
\[
 \Delta_xT(x,y)=T(x+1,y)-T(x,y),\qquad
 \Delta_yT(x,y)=T(x,y+1)-T(x,y),
\]
donde las coordenadas se toman módulo nueve. Esta convención fija también la
orientación positiva de cada plaqueta: primero se recorre la dirección
\(x\), después la dirección \(y\), y se regresa por las aristas opuestas.
Es importante tipar correctamente estos dos objetos. \(S\) y \(P\) son
potenciales de vértice; las variables de enlace son sus diferencias
discretas
\[
A_x^T=\Delta_xT,
\qquad
A_y^T=\Delta_yT.
\]
Para una mezcla ordenada de potenciales \((T_x,T_y)\), definimos la
curvatura de plaqueta por
\[
F(T_x,T_y)
=\Delta_x A_y^{T_y}-\Delta_y A_x^{T_x}
=\Delta_x\Delta_yT_y-\Delta_y\Delta_xT_x.
\]
Como
\[
\Delta_xS=\Delta_yS=1,
\qquad
\Delta_xP=y,
\qquad
\Delta_yP=x,
\]
se obtiene en las \(81\) plaquetas:
\[
F(S,S)=F(P,P)=0,
\qquad
F(S,P)=+1,
\qquad
F(P,S)=-1\pmod9.
\]
Las conexiones asociadas a un único observable son planas; la composición
mixta de los dos observables produce curvatura de signo opuesto según el
orden. Para un rectángulo coordenado \(R_{a,b}\) contenido
en la carta, definimos el Wilson aditivo de borde como la suma orientada de
las variables de enlace. La cancelación de las aristas interiores da
\[
 W_{\partial R_{a,b}}(S,P)=ab,\qquad
 W_{\partial R_{a,b}}(P,S)=-ab\pmod9,
\]
que es la identidad de Stokes finita para esta convención.

La distinción entre potencial y enlace es algebraicamente necesaria. Si se sustituyesen
directamente \(A_x,A_y\) por \(S,P\), aparecerían términos dependientes de
\(x,y\) y no se obtendrían los valores constantes anteriores. La formulación
tipada es la que produce el cociclo mixto utilizado más adelante en el cierre
de Heisenberg finito.

\begin{remark}
La diferencia integrada \(54\) y la curvatura mixta \(\pm1\) son invariantes
distintos. La primera compara sumas globales; la segunda depende del orden de
los dos potenciales y se localiza en cada plaqueta. La formulación mediante
operadores autoadjuntos y conmutadores se desarrollará después de construir el
espacio de estados correspondiente.
\end{remark}

\section{Expansión del operador de transición por caminos}

La definición de APP incluye el espacio
\(\operatorname{Path}(\GraphAPP)\). Una vez fijado un operador de
transición \(U\) sobre los vértices de \(\GraphAPP\), suponemos que
\(U_{yx}=0\) siempre que \(x\to y\) no sea una arista del grafo. Bajo esta
hipótesis de soporte, sus potencias admiten la expansión
\begin{proposition}[Expansión finita por caminos]
\label{prop:app-suma-caminos}
Para cada \(R\ge1\) y cada par de vértices \(i,f\),
\[
(U^R)_{fi}
=\sum_{\substack{\gamma:i\to f\\|\gamma|=R}}
\prod_{k=0}^{R-1}U_{x_{k+1},x_k},
\qquad
\gamma=(x_0=i,x_1,\ldots,x_R=f).
\]
\end{proposition}

\begin{proof}
Para \(R=1\) la identidad es la definición de los elementos de matriz de
\(U\). Si vale para \(R\), entonces
\[
 (U^{R+1})_{fi}
 =\sum_j U_{fj}(U^R)_{ji}.
\]
Sustituir la hipótesis inductiva concatena a cada camino de longitud \(R\)
desde \(i\) hasta \(j\) la última arista \(j\to f\); así, cada camino de
longitud~\(R+1\) aparece exactamente una vez.\qedhere
\end{proof}

En consecuencia, APP reúne en una sola estructura el toro aritmético finito,
los observables de suma y producto, la descomposición residuo--cociente, el
radical nilpotente y una curvatura mixta orientada. El representante \(9\) de
la clase nula fija una carta visible del cociente; reemplazarlo por \(0\) sin
transportar simultáneamente la carta altera los registros empleados por la
dinámica posterior.
\fi
